% © 2025 Ing. David Jaroš — CC BY-NC-ND 4.0
%
% This work is licensed under a Creative Commons Attribution-NonCommercial-NoDerivatives 
% 4.0 International License (CC BY-NC-ND 4.0).
%
% License History: Earlier drafts (up to v0.3) were released under CC BY 4.0. 
% From v0.4 onward, all material is released under CC BY-NC-ND 4.0 to protect 
% the integrity of the theoretical work during ongoing academic development.
%
% See LICENSE.md for full license text.

% This file can be compiled standalone or included in another document
\ifdefined\INCLUDEMODE
  % Being included in another document - skip preamble
\else
  % Standalone compilation - include preamble
  \documentclass[12pt]{article}
  \usepackage[a4paper, margin=2.5cm]{geometry}
  \usepackage{amsmath, amssymb, amsthm}
  \usepackage{hyperref}
  \usepackage{xcolor}
  \usepackage{mdframed}

  \newtheorem{theorem}{Theorem}[section]
  \newtheorem{lemma}[theorem]{Lemma}
  \newtheorem{proposition}[theorem]{Proposition}
  \theoremstyle{definition}
  \newtheorem{definition}[theorem]{Definition}
  \theoremstyle{remark}
  \newtheorem{remark}[theorem]{Remark}

  \newmdenv[backgroundcolor=yellow!20, linecolor=orange, linewidth=1.5pt]{gapbox}
  \newmdenv[backgroundcolor=green!15, linecolor=green!60!black, linewidth=1.5pt]{resultbox}

  \title{\textbf{Step 1: U(1)$_\mathrm{EM}$ Protection at $\varphi = \mathrm{const}$}}
  \author{David Jaroš}
  \date{2026-03-06}

  \begin{document}
  \maketitle

  \begin{abstract}
  We prove that the electromagnetic gauge symmetry $U(1)_\mathrm{EM}$ is \emph{not} broken
  by a constant scalar background $\varphi = \mathrm{const} \neq 0$ in UBT.
  The key observation is that $\varphi$ is the real part of the compact $\psi$-component
  of the biquaternionic field $\Theta$, and it carries zero $U(1)_\mathrm{EM}$ charge.
  Consequently, no photon mass term $m_\gamma^2 = g^2\varphi^2$ is generated.
  \textbf{Verdict: PROVED.}
  \end{abstract}
\fi

\section{Step 1: U(1)$_\mathrm{EM}$ Protection at $\varphi = \mathrm{const}$}
\label{sec:u1_protection}

\textbf{Task reference:} UBT\_v29\_task7, step1\_u1\_protection.
\textbf{Date:} 2026-03-06.
\textbf{Verdict: PROVED.}

\subsection{Setup: Field Content and Charge Assignments}

In UBT the fundamental field is the biquaternionic scalar
\begin{equation}
  \Theta(q,\tau) \in \mathbb{C}\otimes\mathbb{H},
  \qquad \tau = t + i\psi,
\end{equation}
where $t$ is real (physical) time and $\psi \in \mathbb{R}$ is the imaginary time component
(compact circle of radius $R_\psi$, see AXIOM B).

The scalar background $\varphi$ is identified with the vacuum expectation value of the
\emph{real part of the $\psi$-phase component} of $\Theta$:
\begin{equation}
  \varphi \;=\; \langle \operatorname{Re}[\Theta\big|_{\psi\text{-projection}}] \rangle.
\end{equation}
Concretely, expanding $\Theta$ in the $\psi$-circle modes,
\begin{equation}
  \Theta(q,\psi) = \sum_{n\in\mathbb{Z}} \theta_n(q)\, e^{in\psi/R_\psi},
\end{equation}
the constant background $\varphi = \operatorname{const}$ corresponds to the $n=0$ (zero-mode)
component $\theta_0(q) = v \in \mathbb{R}$ being turned on.

\subsection{U(1)$_\mathrm{EM}$ Charge of $\varphi$}

$U(1)_\mathrm{EM}$ acts on $\Theta$ as a phase rotation in the \emph{spatial/spinorial}
sector (see \texttt{canonical/interactions/qed.tex}):
\begin{equation}
  \Theta \;\longrightarrow\; e^{i\alpha(x)}\,\Theta,
  \qquad A_\mu \;\longrightarrow\; A_\mu - \tfrac{1}{e}\partial_\mu\alpha.
\end{equation}
This acts on the $\mathbb{C}$-factor of $\mathbb{C}\otimes\mathbb{H}$, specifically on the
charged spinorial components of $\Theta$.

The $\psi$-zero-mode $\theta_0$ projects onto the \emph{gravity/phase sector} of $\Theta$,
which is the $\mathbb{H}$-factor (real quaternionic part).  This sector is \emph{neutral}
under $U(1)_\mathrm{EM}$ because:
\begin{enumerate}
  \item The quaternionic imaginary units $\{i_1,i_2,i_3\}$ commute with the $U(1)_\mathrm{EM}$
        phase (which acts on the $\mathbb{C}$-factor only).
  \item The $\psi$-circle parameterises a direction in $\tau$-space that is orthogonal to
        the $U(1)_\mathrm{EM}$ fibre; its zero mode $v$ is a real number invariant under
        phase rotations.
\end{enumerate}

\begin{theorem}[U(1)$_\mathrm{EM}$ neutrality of $\varphi$]
\label{thm:phi_neutral}
  The scalar $\varphi = \langle\theta_0\rangle \in \mathbb{R}$ has zero $U(1)_\mathrm{EM}$
  charge:
  \begin{equation}
    e^{i\alpha}\,\varphi\,e^{-i\alpha} = \varphi
    \quad\text{for all }\alpha\in\mathbb{R}.
  \end{equation}
\end{theorem}

\begin{proof}
  Since $\varphi\in\mathbb{R}$ is a real number, and $U(1)_\mathrm{EM}$ acts as multiplication
  by a complex phase on the $\mathbb{C}$-factor, which acts trivially on real scalars, we
  have $e^{i\alpha}\varphi = \varphi\,e^{i\alpha}$ in the algebra, i.e., $\varphi$ commutes
  with the phase.  Being also invariant in absolute value, $\varphi$ is uncharged.
\end{proof}

\subsection{No Photon Mass Term}

In the standard Higgs mechanism, a charged scalar field $\phi$ with charge $q$ develops a
VEV $\langle\phi\rangle = v$, giving the gauge boson a mass
$m^2 = (qev)^2$ from the covariant-kinetic term $|D_\mu\phi|^2$.

For the UBT scalar background $\varphi$, the covariant derivative is
\begin{equation}
  D_\mu\varphi = (\partial_\mu + i\,q_\varphi\, e\, A_\mu)\,\varphi.
\end{equation}
Since $q_\varphi = 0$ (Theorem~\ref{thm:phi_neutral}),
\begin{equation}
  D_\mu\varphi = \partial_\mu\varphi = 0
  \qquad\text{(for }\varphi=\mathrm{const}\text{)}.
\end{equation}
Hence the kinetic term $|D_\mu\varphi|^2$ vanishes identically for constant $\varphi$, and
\begin{equation}
  \boxed{m_\gamma^2 = 0.}
\end{equation}

\begin{resultbox}
\textbf{Conclusion (Step 1, PROVED):}
A constant scalar background $\varphi = \mathrm{const} \neq 0$ does not break $U(1)_\mathrm{EM}$.
The photon remains massless.  No Higgs-like mechanism is operative in the EM sector.
This holds because $\varphi$ lives in the neutral (gravity/phase) sector of the
biquaternionic algebra $\mathbb{C}\otimes\mathbb{H}$.
\end{resultbox}

\subsection{Invariance of $S[\Theta]$}

For completeness we verify directly that the UBT action $S[\Theta]$ is invariant under
$\Theta \to e^{i\alpha(x)}\Theta$ when $\varphi = \mathrm{const}$, i.e., that the
$\varphi$-background does not introduce a symmetry-breaking potential in the EM sector.

The action splits as $S = S_\mathrm{kin} + S_\mathrm{gauge} + S_\mathrm{int}$ (see
Appendix AA for the full action).  The $\varphi$-background enters only through
$S_\mathrm{kin}[\Theta_0]$ where $\Theta_0 = v$ is the zero mode.  Since
$S_\mathrm{kin}[\Theta_0] = \mathrm{const}$ for $\Theta_0 = v \in \mathbb{R}$ and
$U(1)_\mathrm{EM}$ leaves real numbers invariant, the full action satisfies
$\delta_{U(1)_\mathrm{EM}} S = 0$.

\subsection{Summary}

\begin{center}
\begin{tabular}{ll}
\hline
Claim & $U(1)_\mathrm{EM}$ unbroken at $\varphi = \mathrm{const}$ \\
Status & \textbf{PROVED} \\
Charge of $\varphi$ & $q_\varphi = 0$ (real, gravity/phase sector) \\
Photon mass & $m_\gamma^2 = 0$ (exact) \\
Relies on & Algebra structure $\mathbb{C}\otimes\mathbb{H}$, AXIOM B \\
Layer & [L0] — pure biquaternionic geometry \\
\hline
\end{tabular}
\end{center}

\ifdefined\INCLUDEMODE\else
  \end{document}
\fi
